Энкодеры --- это датчики, которые формируют цифровой сигнал в ответ на движение: поворотные энкодеры реагируют на вращение вокруг заданной оси, линейные энкодеры реагируют на движение вдоль рабочего направления.

Существует два основных типа энкодеров: инкрементальные и абсолютные. Инкрементальные датчики в ответ на движение формируют серию импульсов, которая может использоваться для измерения скорости или подаваться на вход счетчика для отслеживания положения. По другому принципу работают абсолютные энкодеры: генерируются битовые последовательности, которые кодируют непосредственно положение.

Основным преимуществом абсолютных датчиков является то, что в случае прерывания или отключения питания устройства датчик может продолжать использоваться для непрерывной работы, информация о положении не теряется. В свою очередь, конструкция инкрементального энкодера не обладает таким свойством, и, поэтому, его необходимо приводить в исходное положение, выполняя <<слепой>> поворот до опорного положения --- нулевой точки (англ. \textit{zero position}). Это важное преимущество абсолютного энкодера, поскольку необходимость инициализации путем поворота до нулевой точки, в общем случае, небезопасна, а также приводит к дополнительным временным затратам.

Инкрементальные энкодеры целесообразно использовать для измерения скорости и относительного перемещения там, где допустим цикл калибровки при включении и не требуется знать абсолютное положение. Абсолютные энкодеры используются в системах, где критично сохранять координату в любой момент времени, скорость возможно получать как производную от измеренного абсолютного положения без установки дополнительных датчиков.

Наблюдается тенденция по увеличению доли устройств, использующих абсолютные энкодеры. Для этого есть несколько ключевых причин: разница в цене по сравнению с инкрементальными энкодерами сократилась, что привело к тому, что абсолютные датчики начали активно применяться не только для критически важных задач.

Другие причины носят более отраслевой характер и связаны с технологическими изменениями: одной из них является возросшая потребность в накоплении информации, поскольку предприятия все чаще подключаются к программным системам, таким как ERP (планирование ресурсов предприятия) и SPC (статистический контроль процессов), что увеличивает потребность в обмене данными между машинами, поскольку все больше устройств и компонентов подключается к полевым шинам при модернизации оборудования; другой является желание и потребность в получении большего количества диагностической информации от оборудования для повышения эффективности профилактического обслуживания и предотвращения непредвиденных отказов и простоев.

Основным функциональным элементом абсолютного энкодера является \textit{измерительная шкала} (кодовый диск или кольцо, например), содержащая последовательность меток, и \textit{считывающее устройство} (считывающая головка), преобразующее данные со шкалы в цифровой код. По способу формирования сигнала и физическом явлении, лежащем в основе построения измерительной меры и конструкции, энкодеры подразделяются на несколько основных групп: механические, резистивные, магнитные, оптические, индуктивные и емкостные. Энкодеры также подразделяются на однооборотные и многооборотные, с ограниченным и неограниченным диапазоном. Каждый тип имеет свои преимущества, ограничения и область применения.  

\subsection{Обзор конструкций}
Рассмотрим основные применяемые на практике конструкции:
\begin{enumerate}

\item \textbf{Емкостные}

Измерение углового положения в конструкциях емкостных энкодеров осуществляется за счет изменения электрической емкости между неподвижными электродами (статором) и подвижными электродами (ротором). Угловое положение вычисляется по изменению емкости.

\textbf{Достоинства}:
\begin{itemize}
    \item Относительная простота изготовления и низкая стоимость производства.
    \item Низкое энергопотребление.
    \item Высокая устойчивость к механическим воздействиям.
\end{itemize}

\textbf{Недостатки}:
\begin{itemize}
    \item Относительно большая погрешность.
    \item Подверженность влиянию пыли и грязи, подвержен влиянию внешних радиопомех, критична соосность элементов
\end{itemize}

В астрономических установках схемы емкостных энкодеров используются редко, в основном в качестве вспомогательных датчиков.  

\item \textbf{Индуктивные}  

Индуктивные системы используют изменение индуктивности катушки или взаимной индуктивности между несколькими катушками при вращении ротора --- в статоре и роторе размещаются проводящие дорожки или катушки, взаимодействие которых изменяется при вращении. К этому типу относятся резольверы (resolvers) и датчики на основе вихревых токов (eddy current).

Современные реализации, например Zettlex IncOder (Celera Motion), а также энкодеры LENZ выполняются в виде печатных плат, что делает их легкими, компактными и устойчивыми к загрязнению.

\textbf{Достоинства}:
\begin{itemize}
    \item Не требуется специальная защита от пыли (немагнитной), влаги, смазки, не подвержены влиянию радиопомех
    \item Высокая стойкость к механическим воздействиям
    \item Низкое энергопотребление
\end{itemize}

\textbf{Недостатки}:
\begin{itemize}
    \item Есть существенные ограничения для датчиков на основе вихревых токов: сильная зависимость от температуры, чувствительность сильно зависит от расстояния между частями энкодера, разрешающая способность зависит от глубины скин-слоя материала статора.
\end{itemize}

\item \textbf{Оптические} 

Оптический принцип является наиболее распространенным в прецизионных системах. Он основан на считывании светового потока, проходящего через шкалу с нанесенным кодом. Наиболее типичная конструкция включает источник света (светодиод или лазер), кодовый диск с прозрачными и непрозрачными участками и фотоприемник. В многодорожечных системах каждая дорожка кодирует один бит абсолютного кода, что требует высокой точности изготовления и юстировки. Современные тенденции развития связаны с переходом на однодорожечные системы, в которых используется последовательность уникальных меток.  

Компании Renishaw и Heidenhain являются лидерами в области оптических абсолютных энкодеров. Система Renishaw RESOLUTE использует однотрековую шкалу с оптическим разрешением до 27 бит и обеспечивает точность до 1 угловой секунды при высокой устойчивости к загрязнениям и вибрациям. Аналогичные решения Heidenhain (серии RCN, ECA, ECN) характеризуются наличием кольцевых измерительных элементов с большим внутренним диаметром, что облегчает их интеграцию в оси вращения телескопов. При этом остаются строгие требования к чистоте оптического канала и стабильности зазора между считывающей головкой и шкалой.  

\textbf{Достоинства}
\begin{itemize}
    \item Наивысшая точность и разрешение среди всех типов энкодеров
    \item Высокая устойчивость к электромагнитным помехам
    \item Возможность высокочастотной работы
\end{itemize}

\textbf{Недостатки}
\begin{itemize}
    \item Чувствительны к любым загрязнениям и температуре, необходима специальная защита
    \item Высокая сложность изготовления и стоимость высокоточных решений
    \item Ограниченный срок службы из-за деградирования светодиодов
\end{itemize}

\item \textbf{Магнитные} 

Магнитные энкодеры работают на основе регистрации магнитного поля, создаваемого намагниченным диском или кольцом. Чувствительный элемент (датчик Холла, магниторезистор или TMR-сенсор) определяет направление и амплитуду поля, что позволяет вычислить угловое положение. 

\textbf{Достоинства}:
\begin{itemize}
    \item Высокое разрешение и чувствительность
    \item Наивысшая устойчивость к внешним (немагнитным) загрязнениям и вибрациям
    \item Низкое энергопотребление
    \item Относительно высокая долговечность
\end{itemize}

\textbf{Недостатки}:
\begin{itemize}
    \item Возможен температурный дрейф, свойства магнита и сенсора меняются с температурой.
    \item Требуют защиты от магнитных помех.
    \item Есть ограничения на применимость --- большинство конструкций не подходит для задач со сквозным валом
    \item Нелинейность выходного сигнала, используются коррекционные алгоритмы
\end{itemize}


\item \textbf{Многооборотные системы.}  

Для систем, где требуется отслеживание положения на угловом диапазоне, превышающем один оборот, применяются многооборотные энкодеры. Традиционные решения используют механические передачи и дополнительные кодовые диски, однако они подвержены износу и накоплению ошибок.  

Современные конструкции применяют электронные методы подсчета числа оборотов. Наибольшее распространение получила технология на основе эффекта Виганда, при котором в специальном проводнике при изменении направления магнитного поля возникает импульс, способный питать счетчик оборотов без внешнего источника. Такая система энергонезависима, не требует батареи и надежно фиксирует число оборотов даже при внезапном отключении питания. Технология используется в многооборотных энкодерах компаний POSITAL и Heidenhain.  
\end{enumerate}

Сравнительный анализ показывает, что оптические энкодеры обеспечивают наилучшие метрологические характеристики: высокое разрешение, низкий уровень шумов и минимальную ошибку, так что они оптимальны для главных осей телескопов. Индуктивные энкодеры представляют собой надежное и менее требовательное решение, обеспечивающее достаточную точность для крупных механических систем. Магнитные энкодеры пригодны для вторичных осей и систем грубого позиционирования, где важна устойчивость к среде. Емкостные решения применяются редко из-за ограниченной точности.  

При разработке прототипа абсолютного датчика угла для астрономического телескопа в этой работе используется оптический принцип измерения. Комбинация источника излучения и линейного пзс-фотоприемника позволяет реализовать схему абсолютного определения углового положения с возможностью цифровой обработки изображения и последующей коррекции ошибок. Такая конструкция может стать технологически доступной альтернативой промышленным высокоточным энкодерам, обеспечивая при этом требуемую точность и стабильность для задач астрономического наблюдения. 

\subsection{Интерфейсы промышленных энкодеров}

\subsubsection{Классификация интерфейсов}

Выходной интерфейс энкодера определяет способ передачи данных об угловом положении в систему управления. По архитектуре передачи данных интерфейсы энкодеров делятся на три категории:
\begin{itemize}
\item \textbf{Параллельные} --- каждый бит выходного кода передается по отдельной паре проводников. Обеспечивают максимальное быстродействие, однако требуют большого числа соединительных линий (по одной паре на каждый бит разрядности), что делает их практически неприменимыми для высокоразрядных абсолютных энкодеров.

\item \textbf{Последовательные} --- биты передаются поочередно по одной или двум парам проводников; являются наиболее распространенным решением для абсолютных датчиков.

\item \textbf{Полевые шины и промышленный Ethernet} (CANopen, PROFIBUS, EtherCAT и др.) — применяются при необходимости интеграции энкодера в распределенную систему управления.
\end{itemize}

\subsubsection{Используемые интерфейсы}

В системах точного позиционирования, типичными интерфейсами для абсолютных энкодеров являются SSI, BiSS‑C и EnDat, а для инкрементных — квадратурный ABZ. Рассмотрим каждый из них подробнее:

\begin{enumerate}
\item Квадратурный интерфейс \textbf{ABZ} является стандартным выходом инкрементальных энкодеров. Он включает три сигнала:
\begin{itemize}
\item \textbf{A} и \textbf{B} --- два идентичных сигнала, сдвинутых друг относительно друга на четверть по фазе (квадратура). По взаимному опережению фаз определяется направление вращения; подсчет фронтов обоих сигналов позволяет реализовать умножение разрешения в 2 или 4 раза.
\item \textbf{Z} (Index) --- одиночный импульс за оборот, используемый как опорная нулевая точка для инициализации счетчика положения
\end{itemize}

Данный интерфейс применяется в прототипе, разрабатываемом в рамках настоящей работы, на первом этапе отработки считывания сигнала с линейки фотоприемников.

\item \textbf{SSI} (Synchronous Serial Interface; синхронный последовательный интерфейс), является простым и широко распространенным стандартом передачи данных абсолютных энкодеров.

Принцип работы: контроллер (ведущее устройство, master) инициирует передачу, генерируя серию тактовых импульсов; энкодер по первому тактовому фронту фиксирует текущее положение в выходном регистре и побитно передает данные синхронно с тактом.

Простота и открытость SSI обусловили его широкое применение в системах точного позиционирования. Недостаток --- отсутствие возможности записи конфигурационных параметров и диагностики энкодера.

\item \textbf{BiSS-C} (Bidirectional Synchronous Serial Interface, version C; открытый последовательный протокол), разработанный как развитие SSI с поддержкой двунаправленной связи.

Ключевые особенности:
\begin{itemize}
\item Высокая тактовая частота, на порядок выше, чем у SSI.
\item Двунаправленная связь: контроллер может и читать положение, и записывать конфигурационные параметры, управлять энкодером и получать диагностическую информацию.
\item Открытый протокол
\item Поддержка топологии <<шина>> --- несколько энкодеров на одной линии.
\end{itemize}

\item \textbf{EnDat} --- проприетарный последовательный протокол компании Heidenhain, применяемый в высокоточных энкодерах этого производителя. По функциональности близок к BiSS-C: поддерживает двунаправленную связь, запись параметров, диагностику и передачу данных.

Применение EnDat в системах управления телескопами обусловлено прежде всего использованием энкодеров Heidenhain, упомянутых в разделе литературного обзора настоящей работы.
\end{enumerate}

\subsubsection*{Сравнительный анализ}

\begin{table}[ht!]
    \centering
    \begin{tabular}{|l|c|c|c|c|}
        \hline
        Параметр & ABZ & SSI & BiSS-C & EnDat 2.2 \\
        \hline
        Тип датчика &
        инк. &
        абс. &
        абс. &
        абс. \\
        \hline
        Напр. связи &
        одно &
        одно &
        два &
        два \\
        \hline
        Физ. ур. &
        RS-422 &
        RS-422 &
        RS-422 &
        RS-422 \\
        \hline
        Макс. частота &
        --- &
        до 1.5 МГц &
        до 10 МГц &
        до 16 МГц \\
        \hline
        Лицензирование &
        открытый &
        открытый &
        открытый &
        проприет. \\
        \hline
    \end{tabular}
    \caption{Сравнение интерфейсов энкодеров}
\end{table}

\subsection{Характеристики энкодеров}

\subsubsection{Основные метрологические параметры}

\textbf{Разрешение} --- это наименьшее угловое перемещение, которое энкодер способен зафиксировать и отобразить в выходном сигнале.

Для инкрементального энкодера разрешение выражается числом импульсов на оборот (PPR, pulses per revolution) либо, при квадратурном подсчете четырех фронтов, числом отсчетов на оборот ($\text{CPR} = 4 \cdot\text{PPR}$).

Для абсолютных энкодеров шаг квантования определяется разрядностью выходного кода: $N$-разрядный энкодер дискретизирует полный оборот $360^\circ$ на $2^N$ уникальных кодовых состояний. Например, 13‑разрядный энкодер обеспечивает $2^{13} = 8192$ дискретных положения на оборот, что соответствует угловому шагу порядка $2.6$ угловой минуты. 

Вместе с тем, при использовании алгоритмов субпиксельной обработки аналоговых выходных сигналов (например, интерполяции синусно-косинусных каналов для инкрементальных энкодеров) эффективное разрешение системы энкодера вместе с его электронным трактом обработки сигнала может существенно превышать формально задаваемое количеством разрядов выходного кода.

Следует подчеркнуть, что разрешение характеризует лишь дискретность представления углового положения и не определяет само по себе точность измерения, то есть степень соответствия измеренного положения действительному углу, ограничиваемую систематическими и случайными погрешностями.

\textbf{Точность} --- это максимальное отклонение измеренного энкодером углового положения от истинного. Точность является интегральной характеристикой, складывающейся из множества источников ошибок: погрешности нанесения шкалы, ошибок считывающей головки, ошибок установки, деформаций вала и эксцентриситета шкалы относительно оси вращения.

Увеличение разрешения не компенсирует систематическую ошибку: энкодер с высоким разрешением, но некачественной шкалой будет повторяемо указывать одно и то же неверное положение.

\textbf{Повторяемость} --- это максимальное расхождение между показаниями энкодера при многократном возврате в одно и то же фактическое положение. Различают однонаправленную повторяемость (подход к заданному положению с одной стороны) и двунаправленную повторяемость (с двух сторон, в том числе для учета люфта приводов).

В задачах астрономических наблюдений, где телескоп многократно наводится на один и тот же объект, высокая повторяемость нередко является более критичным параметром, чем абсолютная точность, поскольку систематическую составляющую ошибки можно компенсировать калибровкой, тогда как случайный разброс устранению не поддается.

\subsection{Кодовые растры}

Кодовый растр оптического энкодера представляет собой систему из одной или нескольких концентрических дорожек с чередованием прозрачных и непрозрачных участков, считываемых головками энкодера. Каждая дорожка формирует один бит в многодорожечных абсолютных энкодерах или один/несколько квадратурных сигналов в инкрементных энкодерах. В общем случае различают два принципиальных подхода к построению оптического растра: многодорожечные (традиционные) и однодорожечные кодовые диски (компактные энкодеры) (см. рис. \ref{fig:track_disk}).

\begin{figure}[ht!]
    \centering
    \includegraphics[width=\linewidth]{BespyatyyIV/images/многодорожечные диски.png}
    \caption{Примеры многодорожечных кодовых дисков}
    \label{fig:track_disk}
\end{figure}

\subsubsection{Многодорожечные абсолютные энкодеры и код Грея.}
В классической реализации абсолютного энкодера используется несколько концентрических дорожек, каждая из которых задает один бит кодового слова. Для углового положения, разделенного на $2^n$ дискретов, формируется $n$ дорожек, и в любой фиксированной угловой позиции диск задает уникальное $n$-битное слово.

Наиболее распространенной формой кодирования в таких системах является циклический код Грея, обеспечивающий изменение только одного бита при переходе между соседними дискретами положения. Это снижает вероятность ошибок считывания при частичном перекрытии фронтов (несовпадении момента переключения всех дорожек) и повышает помехоустойчивость измерения (см. рис. \ref{fig:gray_code_tracks}).

В многодорожечных многооборотных энкодерах кодовый растр дополняется дополнительными дорожками или отдельным кодовым диском, что позволяет измерять как положение в пределах одного оборота, так и число совершенных оборотов.

\begin{figure}[ht!]
    \centering
    \includegraphics[width=0.7\linewidth]{BespyatyyIV/images/двоичный и грей.png}
    \caption{Сравнение для $n=5$}
    \label{fig:gray_code_tracks}
\end{figure}

\subsubsection{Однодорожечные абсолютные энкодеры и псевдослучайные последовательности.}
Для повышения компактности применяются однодорожечные кодовые диски, в которых информация о положении кодируется в виде длинной единой бинарной последовательности. В этом случае вдоль окружности формируется циклическая последовательность, обладающая свойством уникальности всех подслов фиксированной длины.

Типичный подход использует последовательности де Брейна или другие псевдослучайные бинарные коды, для которых каждая возможная комбинация длины $n$ встречается на диске ровно один раз. Система оптических датчиков считывает одновременно $n$ соседних участков дорожки, образуя текущий $n$-битный фрагмент; по этому фрагменту однозначно восстанавливается угловое положение в пределах полного оборота.

Для обеспечения синхронизации и определения направления вращения однодорожечные кодовые диски часто дополняются синхронизационной дорожкой или специализированным маркером (репером).

\subsubsection{Инкрементные энкодеры и квадратурные последовательности.}
Для инкрементных энкодеров основной задачей является формирование равномерной последовательности импульсов, соответствующей приращению угла, без прямого кодирования абсолютного положения. Наиболее распространенный вариант использует две дорожки (или два смещенных по фазе оптических канала) с сигналами A и B, формирующими квадратурный код (фазовый сдвиг $90^\circ$) (см. рис. \ref{fig:quadrature_signals}).

При равномерном вращении кодовый растр формирует две квадратурные последовательности, по которым измерительная электроника определяет число импульсов (угловое смещение) и направление вращения (по относительной фазе сигналов A и B). Как правило, добавляется отдельная референтная дорожка (сигнал Z) для генерации одного импульса на оборот, используемого для привязки счетчика к фиксированному угловому положению.

Разрешение инкрементного энкодера определяется числом периодов сигнала на оборот (числом штрихов), при этом при учете всех фронтов сигналов A и B достигается четырехкратное повышение эффективного числа отсчетов (режим $\times 4$-квадратурного счета). Возможны также многодорожечные инкрементные структуры, в которых на дополнительный канал выносятся, например, сдвинутые по фазе последовательности или сигналы с другим шагом (верньерное измерение \cite{bahn_vernier}, комбинирование нескольких каналов, аналогично принципу работы штангенциркуля см. рис. \ref{fig:vernier_track});

\begin{figure}[ht!]
    \centering
    \includegraphics[width=0.9\linewidth]{BespyatyyIV/images/ABZ инкрементальный диск .png}
    \caption{Квадратурные сигналы A и B, а также референтный импульс Z инкрементного энкодера.}
    \label{fig:quadrature_signals}
\end{figure}

\begin{figure}[ht!]
    \centering
    \includegraphics[width=0.85\linewidth]{BespyatyyIV/images/верньерные дорожки.png}
    \caption{Пример верньерных дорожек инкрементного диска}
    \label{fig:vernier_track}
\end{figure}

\subsubsection{Специализированные и комбинированные кодовые структуры.}

Существуют схемы кодирования абсолютного угла, основанные на оригинальных принципах построения кодового шаблона. В частности, однодорожечные коды Грея \cite{singletrack_gray} \cite{habr_encoder} реализуют абсолютное кодирование по одной кольцевой дорожке (см. рис. \ref{fig:single_gray}) за счет использования нескольких пространственно разнесенных считывающих головок (см. рис. \ref{fig:gray_heads}); при этом метки на шкале размещаются под строго рассчитанными углами, обеспечивающими однозначное восстановление положения по совокупности сигналов с головок. 

\begin{figure}[ht!]
    \centering
    \begin{subfigure}{0.5\textwidth}
        \centering
        \includegraphics[width=0.65\textwidth]{BespyatyyIV/images/однодор грей код.png}
        \caption{Код}
        \label{fig:single_gray}
    \end{subfigure}%
    \begin{subfigure}{0.5\textwidth}
        \centering
        \includegraphics[width=0.65\textwidth]{BespyatyyIV/images/однодор грей головки.png}
        \caption{Расположение головок}
        \label{fig:gray_heads}
    \end{subfigure}
    \caption{Конфигурация описанная в работе \cite{habr_encoder}}
\end{figure}

Другой пример представляет собой кодовая структура вида «репер–код» \cite{mostkom_patent}, в которой шкала разделена на сегменты с выделенными реперными штрихами (опорными уникальными маркерами), а считывающее устройство определяет положения двух соседних реперов и по расположению кодовых меток между ними однозначно идентифицирует номер сегмента и положение внутри сегмента (см. рис. \ref{fig:mostkom_code}).

\begin{figure}[ht!]
    \centering
    \includegraphics[width=0.75\linewidth]{BespyatyyIV/images/мостком конфиг.png}
    \caption{Конфигурация описанная в работе \cite{mostkom_patent}. Состоит из секторов (1), каждый из которых включает код сектора (2) и репер (3).}
    \label{fig:mostkom_code}
\end{figure}

Наконец, существуют комбинированные многодорожечные диски, сочетающие как инкрементальные, так и абсолютные растры (см. рис. \ref{fig:hybrid_track})

\begin{figure}[ht!]
    \centering
    \includegraphics[width=0.9\linewidth]{BespyatyyIV/images/однодорожечный диск.png}
    \caption{Пример комбинированного диска}
    \label{fig:hybrid_track}
\end{figure}

В результате выбор конкретного типа кодового растра (однодорожечного или многодорожечного) и алгоритма кодирования (код Грея, псевдослучайные последовательности, квадратурные последовательности и их комбинации) определяется требованиями к разрешению, габаритам кодового диска, устойчивости к помехам и сложности считывающей электроники.

